\documentclass[../main.tex]{subfiles}
\begin{document}

\section{Lecture 24 -- Boundary Layers}\label{sec:lec24}
\begin{center}
  \textbf{Date:} June 29, 2026
\end{center}

\subsection*{Overview}
The last lecture of the course studies one of the most important singular
limits in fluid mechanics: \textbf{small viscosity near a solid boundary}.  At
large Reynolds number it is tempting to drop viscosity and use ideal-fluid
potential flow everywhere.  This is almost correct in the bulk, but it fails
exactly where boundary conditions are imposed.  A viscous fluid obeys no-slip,
so all components of the velocity vanish at a solid wall.  An ideal fluid only
obeys no-penetration, so it may slide along the wall.

The mismatch is repaired by a thin \textbf{boundary layer}.  Outside it the
flow is essentially the ideal solution.  Inside it the velocity changes rapidly
in the direction normal to the wall, and the viscous term, although multiplied
by a small viscosity, becomes comparable to inertia because the normal
gradients are large.  The method is therefore an example of matched asymptotic
thinking: solve different approximate equations in different regions, then
match them in their common range of validity.

The concrete model in the lecture is the flow near the front stagnation point
of a cylinder or sphere.  Zooming in on the wall turns the body into a flat
plane, and the outer potential flow becomes the universal stagnation-point
field
\[
  U_e(x)=\omega x,\qquad V_e(y)=-\omega y,
\]
where $\omega$ is a strain-rate scale with dimensions $T^{-1}$ (not the
vorticity).  The boundary-layer equations reduce to a third-order nonlinear ODE
for a single dimensionless function.  This is the Hiemenz stagnation-point
boundary layer.

\medskip\noindent\textbf{Topics:}
\begin{enumerate}
  \item Why high-Reynolds-number potential flow fails at a viscous wall.
  \item Local potential-flow expansion near a wall stagnation point.
  \item Derivation of the Prandtl boundary-layer equations.
  \item Boundary conditions and pressure matching to the outer flow.
  \item Scaling estimate for the boundary-layer thickness.
  \item Similarity reduction of the stagnation-point boundary layer.
  \item Shooting method and the numerical wall-shear coefficient.
  \item Physical meaning of the wall shear and the course summary.
\end{enumerate}

%%─────────────────────────────────────────────────────────────────────────────
\subsection{Why the Ideal Approximation Breaks Near a Wall}\label{sec:lec24:why-ideal-breaks}
%%─────────────────────────────────────────────────────────────────────────────

\subsubsection*{Physical Motivation}
Consider a cylinder of radius $R$ moving through a fluid, or equivalently a
fixed cylinder with fluid flowing past it.  For an incompressible viscous fluid
the governing equations are
\begin{equation}
  \boxed{
  \divg\vel=0,\qquad
  \partial_t\vel+(\vel\cdot\grad)\vel
  =-\frac{1}{\rho}\grad p+\nu\lapl\vel .}
  \label{eq:lec24_incompressible_ns}
\end{equation}
\textit{Significance:} This is the equation whose large-$\reynolds$ limit we
are trying to understand; the subtlety is not the bulk equation but the wall
boundary condition.

The Reynolds number is
\begin{equation}
  \boxed{\reynolds=\frac{UR}{\nu},}
  \label{eq:lec24_reynolds}
\end{equation}
where $U$ is a typical speed, $R$ a typical size, and $\nu=\eta/\rho$ the
kinematic viscosity.  The limit of small viscosity is $\reynolds\gg1$.  A naive
estimate compares
\begin{equation}
  (\vel\cdot\grad)\vel\sim\frac{U^2}{R},
  \qquad
  \nu\lapl\vel\sim\nu\frac{U}{R^2}.
  \label{eq:lec24_bulk_scaling}
\end{equation}
The ratio of viscous to inertial terms is then
\begin{equation}
  \boxed{
  \frac{\nu U/R^2}{U^2/R}\sim\frac{1}{\reynolds}\ll1 .}
  \label{eq:lec24_viscous_bulk_ratio}
\end{equation}
\textit{Significance:} Away from boundaries this justifies the ideal-fluid
approximation at large Reynolds number.

The problem is that the estimate
\[
  \grad\sim 1/R
\]
is not uniformly true.  At a solid wall a viscous fluid obeys
\begin{equation}
  \boxed{\vel=\bvec0\qquad\text{on the wall}.}
  \label{eq:lec24_no_slip}
\end{equation}
\textit{Significance:} No-slip fixes both the normal and tangential components,
so it is stronger than the ideal-fluid no-penetration condition.

An ideal fluid, by contrast, only requires
\begin{equation}
  \boxed{\vel\cdot\uvec n=0\qquad\text{on the wall}.}
  \label{eq:lec24_no_penetration}
\end{equation}
\textit{Significance:} Potential flow may slide tangentially along the wall, so
it cannot satisfy the viscous boundary condition.

Thus, if the outer ideal solution predicts a tangential velocity of order $U$
right next to the wall, the true viscous velocity must fall from $U$ to $0$ over
a short normal distance.  In that thin region
\[
  |\partial_y\vel|\gg U/R,
\]
so the viscous term cannot be dropped.  The small parameter multiplies the
highest spatial derivatives, and the limit $\nu\to0$ is singular.

\begin{figure}[h]
\centering
\begin{tikzpicture}[scale=1.0, >=Stealth]
  \fill[blue!8] (0,0) rectangle (7,0.75);
  \draw[very thick] (0,0) -- (7,0) node[right] {wall};
  \draw[->, thick] (0.6,0) -- (0.6,2.2) node[above] {$y$};
  \draw[->, thick] (0.6,0) -- (2.2,0) node[below] {$x$};
  \draw[blue!70!black, thick] (0,0.75) -- (7,0.75);
  \node[blue!70!black] at (6.0,1.05) {$y\sim\delta$};
  \foreach \x/\len in {1.1/0.35,2.0/0.65,2.9/0.95,3.8/1.25,4.7/1.55,5.6/1.85,6.5/2.15} {
    \draw[->, myred, thick] (\x,1.25) -- ++(\len,0);
  }
  \foreach \x/\len in {1.1/0.08,2.0/0.16,2.9/0.24,3.8/0.32,4.7/0.40,5.6/0.48,6.5/0.56} {
    \draw[->, vcol, thick] (\x,0.35) -- ++(\len,0);
  }
  \node[myred] at (6.0,1.65) {outer flow $U_e(x)$};
  \node[vcol] at (2.8,0.42) {rapid adjustment};
  \node[blue!70!black] at (2.0,0.95) {boundary layer};
\end{tikzpicture}
\caption{A boundary layer is thin in the wall-normal direction.  The tangential
velocity is zero at the wall but matches the outer ideal flow across a distance
$\delta\ll R$.}
\label{fig:lec24_boundary_layer_sketch}
\end{figure}

%%─────────────────────────────────────────────────────────────────────────────
\subsection{Local Potential Flow Near a Stagnation Point}\label{sec:lec24:outer-stagnation}
%%─────────────────────────────────────────────────────────────────────────────

\subsubsection*{Mathematical Setup}
We zoom in on the front stagnation point of the body.  At that scale the curved
surface looks like the plane $y=0$, the fluid occupies $y>0$, and the
stagnation point is at $(x,y)=(0,0)$.  The outer flow is ideal,
incompressible, and irrotational, so
\begin{equation}
  \vel=\grad\phi,\qquad \lapl\phi=0.
  \label{eq:lec24_potential_laplace}
\end{equation}
At the stagnation point the velocity vanishes:
\begin{equation}
  \grad\phi(0,0)=\bvec0.
  \label{eq:lec24_stagnation_condition}
\end{equation}

Expand $\phi$ to second order:
\begin{equation}
  \phi(x,y)=\phi_0+\frac12\left(Ax^2+2Bxy+Cy^2\right)+O(r^3).
  \label{eq:lec24_phi_taylor}
\end{equation}
The constant $\phi_0$ does not affect the velocity, and the linear terms are
absent because of \eqref{eq:lec24_stagnation_condition}.  No-penetration on the
wall requires
\begin{equation}
  v_y(x,0)=\partial_y\phi(x,0)=Bx=0
  \quad\text{for every }x,
  \label{eq:lec24_wall_no_penetration_taylor}
\end{equation}
so $B=0$.  Laplace's equation gives
\begin{equation}
  \lapl\phi=A+C=0,
  \label{eq:lec24_laplace_taylor}
\end{equation}
so $C=-A$.  Writing $A=\omega$, the local potential is
\begin{equation}
  \boxed{\phi(x,y)=\frac{\omega}{2}\left(x^2-y^2\right),\qquad
  \vel_e=(U_e,V_e)=(\omega x,-\omega y).}
  \label{eq:lec24_outer_stagnation_flow}
\end{equation}
\textit{Significance:} This hyperbolic velocity field is the universal leading
form of two-dimensional potential flow near a wall stagnation point.

On the wall $y=0$ this ideal solution gives
\begin{equation}
  U_e(x)=\omega x,\qquad V_e(x,0)=0.
  \label{eq:lec24_outer_wall_velocity}
\end{equation}
The normal boundary condition is satisfied, but the tangential velocity is not
zero except at the stagnation point itself.  This is precisely the mismatch the
boundary layer must correct.

The same outer solution fixes the pressure variation along the layer.  For a
steady incompressible potential flow, Bernoulli gives
\begin{equation}
  p_e+\frac12\rho|\vel_e|^2=\text{const}.
  \label{eq:lec24_outer_bernoulli}
\end{equation}
Near the wall, $|\vel_e|^2\simeq U_e^2=\omega^2x^2$, hence
\begin{equation}
  \boxed{
  p_e(x)=p_s-\frac12\rho\omega^2x^2,\qquad
  -\frac{1}{\rho}\dv{p_e}{x}=\omega^2x=U_e\dv{U_e}{x}.}
  \label{eq:lec24_outer_pressure_gradient}
\end{equation}
\textit{Significance:} The boundary layer does not determine its own pressure
to leading order; it inherits the tangential pressure gradient from the outer
ideal flow.

The complex-potential language from planar potential flow also applies:
\begin{equation}
  w(z)=\phi+i\psi=\frac{\omega}{2}z^2,\qquad z=x+iy,
  \label{eq:lec24_complex_stagnation_potential}
\end{equation}
so
\begin{equation}
  \psi(x,y)=\omega xy.
  \label{eq:lec24_stream_function_outer}
\end{equation}
The streamlines are therefore hyperbolae $xy=\text{const}$.

%%─────────────────────────────────────────────────────────────────────────────
\subsection{Prandtl Boundary-Layer Equations}\label{sec:lec24:prandtl-equations}
%%─────────────────────────────────────────────────────────────────────────────

\subsubsection*{Mathematical Setup}
Inside the boundary layer write the velocity components as
\[
  \vel=(u(x,y),v(x,y)),
\]
where $u$ is tangential to the wall and $v$ is normal to it.  We consider a
steady incompressible flow.  The component form of Navier--Stokes is
\begin{equation}
  \begin{aligned}
  u\partial_xu+v\partial_yu
  &=-\frac1\rho\partial_xp
    +\nu\left(\partial_x^2u+\partial_y^2u\right),\\
  u\partial_xv+v\partial_yv
  &=-\frac1\rho\partial_yp
    +\nu\left(\partial_x^2v+\partial_y^2v\right),\\
  \partial_xu+\partial_yv&=0.
  \end{aligned}
  \label{eq:lec24_component_ns}
\end{equation}

Let $L$ be the tangential scale and $\delta$ the boundary-layer thickness, with
$\delta\ll L$.  The tangential velocity has size $U$.  Continuity gives the
scale of the normal velocity:
\begin{equation}
  \partial_xu\sim\frac{U}{L},\qquad
  \partial_yv\sim\frac{V}{\delta},
  \qquad
  \partial_xu+\partial_yv=0
  \quad\Rightarrow\quad
  \boxed{V\sim\frac{\delta}{L}U\ll U.}
  \label{eq:lec24_normal_velocity_scaling}
\end{equation}
\textit{Significance:} The normal velocity is small, but its normal derivative
is not; this is why it can still enter the tangential momentum equation.

Now compare terms in the tangential equation.  The two advective terms scale as
\begin{equation}
  u\partial_xu\sim\frac{U^2}{L},\qquad
  v\partial_yu\sim
  \left(\frac{\delta}{L}U\right)\frac{U}{\delta}
  =\frac{U^2}{L}.
  \label{eq:lec24_advection_scaling}
\end{equation}
The viscous terms scale as
\begin{equation}
  \nu\partial_x^2u\sim\nu\frac{U}{L^2},\qquad
  \nu\partial_y^2u\sim\nu\frac{U}{\delta^2}.
  \label{eq:lec24_viscous_scaling}
\end{equation}
Since $\delta\ll L$, the $y$-derivative viscous term dominates the $x$-derivative
viscous term.  The leading balance is therefore
\begin{equation}
  \frac{U^2}{L}\sim\nu\frac{U}{\delta^2}.
  \label{eq:lec24_thickness_balance_raw}
\end{equation}
Solving for $\delta$ gives
\begin{equation}
  \boxed{\delta^2\sim\frac{\nu L}{U},\qquad
  \frac{\delta}{L}\sim\reynolds^{-1/2}.}
  \label{eq:lec24_boundary_layer_thickness_general}
\end{equation}
\textit{Significance:} Even as $\nu\to0$, the product
$\nu\,\partial_y^2u$ remains important because
$\partial_y^2u\sim U/\delta^2$ diverges like $1/\nu$.

The normal momentum equation is smaller.  To leading order it says that
pressure does not change across the thin layer:
\begin{equation}
  \boxed{\partial_y p=0,\qquad p(x,y)=p_e(x).}
  \label{eq:lec24_pressure_constant_across_layer}
\end{equation}
\textit{Significance:} Boundary-layer pressure is imposed by the outer inviscid
flow; the layer adjusts velocity and shear, not pressure across its thickness.

Putting the leading terms together gives the Prandtl boundary-layer equations:
\begin{equation}
  \boxed{
  u\partial_xu+v\partial_yu
  =-\frac1\rho\dv{p_e}{x}+\nu\partial_y^2u,\qquad
  \partial_xu+\partial_yv=0,\qquad
  \partial_y p=0.}
  \label{eq:lec24_prandtl_equations}
\end{equation}
\textit{Significance:} These equations retain exactly the viscous term needed
to impose no-slip while discarding terms smaller by powers of $\delta/L$.

Using the outer Euler equation,
\begin{equation}
  -\frac1\rho\dv{p_e}{x}=U_e(x)\dv{U_e}{x},
  \label{eq:lec24_outer_euler_pressure}
\end{equation}
the tangential Prandtl equation may be written as
\begin{equation}
  \boxed{
  u\partial_xu+v\partial_yu
  =U_e\dv{U_e}{x}+\nu\partial_y^2u.}
  \label{eq:lec24_prandtl_with_outer_velocity}
\end{equation}
\textit{Significance:} This is the matching form: the outer ideal flow enters
only through the function $U_e(x)$.

The boundary conditions are
\begin{equation}
  \boxed{
  u(x,0)=0,\qquad v(x,0)=0,\qquad
  u(x,y\to\infty)\to U_e(x).}
  \label{eq:lec24_boundary_layer_bcs}
\end{equation}
\textit{Significance:} The layer satisfies the viscous wall condition at
$y=0$ and matches the ideal solution as one exits the thin region.

%%─────────────────────────────────────────────────────────────────────────────
\subsection{Similarity Solution for Stagnation-Point Flow}\label{sec:lec24:similarity}
%%─────────────────────────────────────────────────────────────────────────────

\subsubsection*{Motivation}
For the stagnation-point outer flow
\[
  U_e(x)=\omega x,
\]
the boundary-layer thickness estimate becomes especially simple:
\begin{equation}
  \boxed{\delta\sim\sqrt{\frac{\nu}{\omega}}.}
  \label{eq:lec24_stagnation_thickness}
\end{equation}
\textit{Significance:} Near a stagnation point the local strain rate $\omega$
sets the tangential gradient, so the layer thickness is independent of $x$ in
the leading similarity approximation.

To solve the equations exactly within this approximation, introduce the
dimensionless normal coordinate
\begin{equation}
  \tilde y=y\sqrt{\frac{\omega}{\nu}},
  \label{eq:lec24_dimensionless_y}
\end{equation}
and a stream function
\begin{equation}
  \boxed{\psi(x,y)=\sqrt{\nu\omega}\,x\,f(\tilde y),}
  \label{eq:lec24_stream_function_ansatz}
\end{equation}
\textit{Significance:} The stream function automatically satisfies
incompressibility, and the chosen prefactor gives the correct dimensions.

With the convention
\begin{equation}
  u=\partial_y\psi,\qquad v=-\partial_x\psi,
  \label{eq:lec24_stream_velocity_definition}
\end{equation}
we obtain
\begin{equation}
  \boxed{
  u(x,y)=\omega x f'(\tilde y),\qquad
  v(x,y)=-\sqrt{\nu\omega}\,f(\tilde y),}
  \label{eq:lec24_similarity_velocities}
\end{equation}
where primes denote derivatives with respect to $\tilde y$.
\textit{Significance:} The whole velocity field is determined by one
dimensionless profile $f$.

\subsubsection*{Derivation of the ODE}
Compute the needed derivatives:
\begin{equation}
  \begin{aligned}
  \partial_xu&=\omega f',\\
  \partial_yu&=\omega x f''\sqrt{\frac{\omega}{\nu}},\\
  \partial_y^2u&=\omega x f'''\frac{\omega}{\nu}.
  \end{aligned}
  \label{eq:lec24_similarity_derivatives}
\end{equation}
Therefore
\begin{equation}
  u\partial_xu
  =(\omega x f')(\omega f')
  =\omega^2x(f')^2,
  \label{eq:lec24_similarity_first_advection}
\end{equation}
and
\begin{equation}
  v\partial_yu
  =\left(-\sqrt{\nu\omega}\,f\right)
   \left(\omega x f''\sqrt{\frac{\omega}{\nu}}\right)
  =-\omega^2x\,f f''.
  \label{eq:lec24_similarity_second_advection}
\end{equation}
The viscous term is
\begin{equation}
  \nu\partial_y^2u=\nu\left(\omega x f'''\frac{\omega}{\nu}\right)
  =\omega^2x f'''.
  \label{eq:lec24_similarity_viscous}
\end{equation}
Since $U_e\,dU_e/dx=(\omega x)\omega=\omega^2x$, substitution into
\eqref{eq:lec24_prandtl_with_outer_velocity} gives
\begin{equation}
  \omega^2x\left[(f')^2-f f''\right]
  =\omega^2x\left[1+f'''\right].
  \label{eq:lec24_similarity_before_dividing}
\end{equation}
For $x\neq0$ divide by $\omega^2x$ and rearrange:
\begin{equation}
  \boxed{f'''+f f''+1-(f')^2=0.}
  \label{eq:lec24_hiemenz_equation}
\end{equation}
\textit{Significance:} The partial differential boundary-layer problem has
collapsed to a nonlinear third-order ordinary differential equation.

The boundary conditions follow from \eqref{eq:lec24_boundary_layer_bcs}.  At
$y=0$,
\[
  u=0\Rightarrow f'(0)=0,\qquad
  v=0\Rightarrow f(0)=0.
\]
Matching $u\to U_e=\omega x$ as $y\to\infty$ gives $f'\to1$.  Thus
\begin{equation}
  \boxed{f(0)=0,\qquad f'(0)=0,\qquad f'(\infty)=1.}
  \label{eq:lec24_hiemenz_bcs}
\end{equation}
\textit{Significance:} The three boundary conditions match the order of the
ODE, but one of them lives at infinity, so the solution is naturally found by
shooting.

\subsubsection*{Shooting and the wall-shear number}
The Hiemenz equation has no elementary closed-form solution.  Numerically one
chooses a trial value of $f''(0)$, integrates the ODE outward from
\[
  f(0)=0,\qquad f'(0)=0,\qquad f''(0)=s,
\]
and adjusts $s$ until the far-field condition $f'(\infty)=1$ is satisfied.
This is the shooting method.  The successful value is
\begin{equation}
  \boxed{f''(0)\simeq 1.2326.}
  \label{eq:lec24_hiemenz_wall_shear_number}
\end{equation}
\textit{Significance:} This single dimensionless number fixes the wall shear in
the stagnation-point boundary layer.

The wall shear stress is the tangential viscous stress
\begin{equation}
  \tau_w(x)=\eta\left.\partial_yu\right|_{y=0},
  \label{eq:lec24_wall_shear_definition}
\end{equation}
where $\eta=\rho\nu$ is the dynamic viscosity.  Using
\eqref{eq:lec24_similarity_derivatives},
\begin{equation}
  \boxed{
  \tau_w(x)
  =\eta\,\omega x\,f''(0)\sqrt{\frac{\omega}{\nu}}
  =\rho f''(0)\,x\,\sqrt{\nu\omega^3}.}
  \label{eq:lec24_wall_shear_result}
\end{equation}
\textit{Significance:} The shear tends to zero as $\nu^{1/2}$, not as $\nu$,
because the velocity gradient grows like $1/\delta\sim\nu^{-1/2}$.

Equivalently, for a general boundary-layer estimate,
\begin{equation}
  \boxed{\tau_w\sim\eta\frac{U}{\delta}
  \sim\rho U\sqrt{\frac{\nu U}{L}}.}
  \label{eq:lec24_shear_scaling_general}
\end{equation}
\textit{Significance:} A very thin layer converts a small viscosity into a
finite leading correction to ideal-flow drag estimates.

%%─────────────────────────────────────────────────────────────────────────────
\subsection{Matched Asymptotic Meaning}\label{sec:lec24:matched-asymptotics}
%%─────────────────────────────────────────────────────────────────────────────

\subsubsection*{Conceptual Background}
The boundary-layer construction is not just a trick for this cylinder example.
It is a general method for singular perturbation problems:
\begin{itemize}
  \item Solve an \textbf{outer problem} in the bulk, where the small term is
  genuinely negligible.
  \item Solve an \textbf{inner problem} near the boundary, using stretched
  variables such as $\tilde y=y/\delta$.
  \item Match the inner solution as $\tilde y\to\infty$ to the outer solution as
  $y\to0$.
\end{itemize}

In the present problem the outer solution is the ideal potential flow.  The
inner solution is governed by Prandtl's equations.  The matching condition is
\[
  u(x,y\to\infty)\to U_e(x).
\]
One may also view the boundary layer as generating an effective boundary
condition for the outer flow: instead of resolving the thin layer everywhere,
one encodes its effect on the bulk through shear stress, displacement
thickness, or separation criteria.

\nt{The lecture only analyzed the neighborhood of the front stagnation point.
Farther along the body, boundary layers can thicken, separate, and feed the
wake behind the object.  That wake is one of the physical ways viscosity
resolves D'Alembert's paradox.}

%%─────────────────────────────────────────────────────────────────────────────
\subsection{Course Perspective}\label{sec:lec24:course-perspective}
%%─────────────────────────────────────────────────────────────────────────────

\subsubsection*{Motivation}
The lecture ended with a short course summary.  The course had two main parts.
The first part, mechanics of solids, introduced stress, strain, Hooke's law,
linear isotropic elasticity, torsion, bending, and the Euler--Bernoulli beam
equation.  The second and larger part, mechanics of fluids, moved through
hydrostatics, ideal-fluid dynamics, Bernoulli's equation, circulation and
vorticity, potential flow, viscosity and Navier--Stokes, exact viscous-flow
solutions, Reynolds-number regimes, sound waves, the action-principle
viewpoint, and finally boundary layers.

The recurring mathematical language was the language of fields:
\[
  \bvec u(\bvec x),\qquad \stress(\bvec x),\qquad \strain(\bvec x),\qquad
  \vel(\bvec x,t),\qquad p(\bvec x,t),\qquad \rho(\bvec x,t).
\]
The recurring physical themes were conservation laws, boundary conditions,
symmetry, dimensional analysis, and the fact that approximate equations have
domains of validity.  Boundary layers are a fitting final topic because they
show all of these themes at once: a conservation law gives Prandtl's equations,
the no-slip boundary condition creates the layer, dimensional analysis predicts
its thickness, and matching tells us how local and global solutions fit
together.

\subsection*{Key Takeaways}
\begin{itemize}
  \item Large $\reynolds$ does not mean viscosity is negligible everywhere; it
  fails near no-slip walls.
  \item Near a wall stagnation point the outer potential flow is
  $\phi=\frac{\omega}{2}(x^2-y^2)$, so $U_e(x)=\omega x$.
  \item In a boundary layer, $\partial_y$ derivatives dominate
  $\partial_x$ derivatives, while pressure is inherited from the outer flow.
  \item The Prandtl equations are the leading high-$\reynolds$ wall equations:
  they keep $\nu\partial_y^2u$ but drop smaller viscous terms.
  \item The general thickness estimate is
  $\delta/L\sim\reynolds^{-1/2}$; for stagnation flow
  $\delta\sim\sqrt{\nu/\omega}$.
  \item The stagnation-point boundary layer reduces to
  $f'''+f f''+1-(f')^2=0$ with
  $f(0)=f'(0)=0$ and $f'(\infty)=1$.
  \item Wall shear scales as $\tau_w\sim\eta U/\delta$, so it vanishes like
  $\nu^{1/2}$ in the small-viscosity limit, not like $\nu$.
\end{itemize}

\end{document}
